\cotefile{63}
\scanpage{9}\mspage{6}
\makeatletter\tagsleft@true\makeatother
5. \uline{Application}: \uline{The} \uline{relation} \uline{between} $L_v(1)$ \uline{and} $\int_{E_v}|\omega_v|$ \uline{in} \uline{case} $k$ \uline{is} \uline{finite}.

Suppose now $R$ is complete and $k$ is finite with $q=\operatorname{Card} k$. On $K$ we agree to use the additive Haar measure with respect to which $R$ has measure $1$. This being agreed, a differential \inserted{$\omega$} on $E$ gives us a measure $|\omega|$ on $E(K)$ in the usual manner.

\uline{Corollary 5.1} (cf theorem 4.2) $\displaystyle\int_{E_1(K)}|\omega|=\frac{1}{q}$, if $\omega$ \uline{is} \uline{a} \uline{differential} \uline{of} $1^{\mathrm{st}}$ \uline{kind} \uline{on} $E$ \uline{coming} \uline{from} \uline{a} \uline{minimal} \uline{equation}.

Indeed, by (4.2) we have $\displaystyle\int_{E_1(K)}|\omega|\unsure{-|=}\int_{\mathfrak p}|d\unsure{z|x}|=\frac{\int_R|dz|}{(R:\mathfrak p)}=\frac{1}{q}$

The local factor \unsure{occuring|occurring} in the Euler product with a good functional equation is (Serre tells me) as follows, in which $N_v=\operatorname{Card}(\widetilde E_0(k))$:
{\small
\[
L_v(s)=\left\{
\begin{array}{@{}c@{\quad}l@{}}
\dfrac{1}{1-(q+1-N_v){\unsure{q|1}}^{-s}+q^{1-2s}}
 & \parbox{.39\linewidth}{, if $\widetilde E$ non singular}\\[1.0ex]
\dfrac{1}{1-q^{-s}}
 & \parbox{.39\linewidth}{, if $\widetilde E$ has a node with the two tangents ratl over $k$ (in which case $N_v=q-1$, and $\widetilde E_0$ is the multiplicative group.)}\\[1.0ex]
\dfrac{1}{1+q^{-s}}
 & \parbox{.39\linewidth}{, if $\widetilde E$ has a node with \unsure{irrational} tangents (in which case $N_v=q+1$, and $\widetilde E_0$ is the twisted mult. group\unsure{.|{}})}\\[1.0ex]
1 & \parbox{.39\linewidth}{, if $\widetilde E$ has a cusp (in which case $N_v=q$, and $\widetilde E_0$ is the additive gp.)}
\end{array}\right.
\tag*{(5.2)}
\]
}

In all cases therefore, $L_v(1)=\dfrac{q}{N_v}$. Since $N_v=(E_0(K):E_1(K))$\unsure{:|{}} it follows from cor 5.1 that we have $\displaystyle\int_{E_0(K)}|\omega|=(L_v(1))^{-1}$, because $|\omega|$ is invariant under translation. Finally then,

\uline{Theorem 5.2}: \uline{If} \uline{we} \uline{use} \uline{the} \uline{measure} \uline{on} $K$ \uline{for} \uline{which} $R$ \uline{gets} \uline{measure} $1$, \uline{and} \uline{use} \uline{a} \uline{differential} \uline{of} \uline{first} \uline{kind} $\omega$ \uline{coming} \uline{from} \uline{a} \uline{minimal} \uline{equation}, \uline{then}
\[
\int_{E(K)}|\omega|=\frac{(E(K):E_0(K))}{L_v(1)}
\tag*{(5.3)}
\]
In other words, the ``fudge factors''\ of Birch and \unsure{Sw-Dyer|Sw.-Dyer} are just the indices $(E(K):E_0(K))$,
% archive sheet number: 8
\clearpage
\scanpage{10}\mspage{7.}
6. \uline{The} \uline{Neron} \uline{minimum} \uline{model}. Suppose $k$ algebraically closed. One can find a regular scheme $X$ over $R$ such that $X\underset{R}{\times}K\approx E$ and such that $X$ is ``minimal'' relative to the map $X\longrightarrow\operatorname{Spec} R$ (i.e such that that map cannot be factored $X\longrightarrow X'\longrightarrow\operatorname{Spec} R$ in such a way that $X\underset{R}{\times}K\to X'\underset{R}{\times}K$ is an isomorphism. Such an $X$ is unique up to isom. Its fiber $\widetilde X=X\underset{R}{\times}k$ is of one of the following types

\diagram{%
\setlength{\arraycolsep}{2.7pt}\renewcommand{\arraystretch}{1.35}%
\resizebox{\linewidth}{!}{$%
\begin{array}{|l|*{10}{c|}}
\hline
\multicolumn{11}{|@{}c@{}|}{%
\begin{tikzpicture}[x=.1726mm,y=.1726mm,line cap=round,baseline=(current bounding box.center)]
\path[use as bounding box] (0,0) rectangle (1450,81);
\begin{scope}
\clip (0,0) rectangle (1450,81);
\foreach \divider in {278,394,510,630}{\draw[line width=.65pt] (\divider,0) -- (\divider,81);}
\node at (140,37) {\deleted{\unsure{Kodaira} symbol}};
\node at (342,34) {$\deleted{\mathrm I_0}$};
\node at (454,34) {$\deleted{\mathrm I_{\unsure{b|\nu}}^{\unsure{\nu|*}}}$};
\node at (575,32) {$\deleted{\mathrm I_{\unsure{b|0}}}$};
\node at (686,33) {$\deleted{\mathrm I_{\unsure{0|b}}^{*}}$};
\foreach \hatchx in {-84,-47,...,1530}{
  \draw[line width=.47pt] (\hatchx,0) -- ++(72,81);
  \draw[line width=.46pt] (\hatchx,81) -- ++(78,-81);}
\end{scope}
\end{tikzpicture}}\\ \hline
\makebox[46mm][l]{Kodaira symbol} &
\makebox[18.2mm]{$\mathrm I_0$} &
\makebox[18.2mm]{$\mathrm I_{\unsure{\nu|v}}\;(\nu>0)$} &
\makebox[18.2mm]{$\mathrm{II}$} &
\makebox[18.2mm]{$\mathrm{III}$} &
\makebox[18.2mm]{$\mathrm{IV}$} &
\makebox[18.2mm]{$\mathrm I_0^{*}$} &
\makebox[18.2mm]{$\mathrm I_\nu^{*}\;(\nu>0)$} &
\makebox[18.2mm]{$\mathrm{IV}^{*}$} &
\makebox[18.2mm]{$\mathrm{III}^{*}$} &
\makebox[18.2mm]{$\mathrm{II}^{*}$} \\ \hline
\shortstack[l]{Picture\\\unsure{.} (the numbers\\indicate multiplicities)} &
\parbox[c][5.45cm][c]{1.82cm}{\figure{%
\begin{tikzpicture}[x=.0158cm,y=-.0158cm,line width=.8pt,line cap=round,line join=round,every node/.style={inner sep=0pt,font=\scriptsize}]
\path[use as bounding box] (0,0) rectangle (112,318);
\draw (51,59) .. controls (77,44) and (79,73) .. (77,130) .. controls (78,189) and (72,233) .. (57,237) .. controls (33,246) and (26,216) .. (29,169) .. controls (30,122) and (38,65) .. (51,59) -- cycle;
\node at (88,146) {$1$};
\end{tikzpicture}}} &
\parbox[c][5.45cm][c]{1.82cm}{\figure{%
\begin{tikzpicture}[x=.0158cm,y=-.0158cm,line width=.8pt,line cap=round,line join=round,every node/.style={inner sep=0pt,font=\scriptsize}]
\path[use as bounding box] (0,0) rectangle (112,318);
\draw (14,164) .. controls (18,130) and (20,88) .. (28,55);
\draw (16,120) .. controls (23,100) and (32,61) .. (43,40);
\draw (28,55) -- (70,48);
\draw (56,39) .. controls (59,60) and (69,81) .. (71,103);
\draw (91,198) -- (88,251);
\draw (61,289) .. controls (76,255) and (82,240) .. (95,218);
\draw (22,305) -- (82,290);
\draw (23,276) -- (42,306);
\node at (30,137) {$1$};
\node at (46,87) {$1$};
\node at (52,63) {$1$};
\node at (76,69) {$1$};
\node at (99,202) {$1$};
\node at (81,267) {$1$};
\node at (55,284) {$1$};
\node at (39,284) {$1$};
\fill (84,143) circle[radius=.45pt];
\fill (83,160) circle[radius=.45pt];
\fill (85,181) circle[radius=.45pt];
\end{tikzpicture}}} &
\parbox[c][5.45cm][c]{1.82cm}{\figure{%
\begin{tikzpicture}[x=.0158cm,y=-.0158cm,line width=.8pt,line cap=round,line join=round,every node/.style={inner sep=0pt,font=\scriptsize}]
\path[use as bounding box] (0,0) rectangle (112,318);
\draw (30,70) .. controls (48,107) and (51,128) .. (53,170) .. controls (52,133) and (53,110) .. (74,69);
\node at (78,90) {$1$};
\end{tikzpicture}}} &
\parbox[c][5.45cm][c]{1.82cm}{\figure{%
\begin{tikzpicture}[x=.0158cm,y=-.0158cm,line width=.8pt,line cap=round,line join=round,every node/.style={inner sep=0pt,font=\scriptsize}]
\path[use as bounding box] (0,0) rectangle (112,318);
\draw (31,76) .. controls (56,135) and (58,175) .. (32,235);
\draw (71,76) .. controls (44,135) and (39,199) .. (75,241);
\node at (28,99) {$1$};
\node at (78,93) {$1$};
\end{tikzpicture}}} &
\parbox[c][5.45cm][c]{1.82cm}{\figure{%
\begin{tikzpicture}[x=.0158cm,y=-.0158cm,line width=.8pt,line cap=round,line join=round,every node/.style={inner sep=0pt,font=\scriptsize}]
\path[use as bounding box] (0,0) rectangle (112,318);
\draw (55,72) -- (60,248);
\draw (23,115) -- (95,195);
\draw (92,108) -- (18,194);
\node at (47,66) {$1$};
\node at (19,107) {$1$};
\node at (91,103) {$1$};
\end{tikzpicture}}} &
\parbox[c][5.45cm][c]{1.82cm}{\figure{%
\begin{tikzpicture}[x=.0158cm,y=-.0158cm,line width=.8pt,line cap=round,line join=round,every node/.style={inner sep=0pt,font=\scriptsize}]
\path[use as bounding box] (0,0) rectangle (112,318);
\draw (62,47) .. controls (57,98) and (70,148) .. (68,283);
\draw (22,38) -- (92,91);
\draw (20,76) -- (91,135);
\draw (35,251) -- (91,190);
\draw (34,288) -- (90,234);
\node at (36,37) {$1$};
\node at (33,65) {$1$};
\node at (77,151) {$2$};
\node at (37,221) {$1$};
\node at (42,264) {$1$};
\end{tikzpicture}}} &
\parbox[c][5.45cm][c]{1.82cm}{\figure{%
\begin{tikzpicture}[x=.0158cm,y=-.0158cm,line width=.8pt,line cap=round,line join=round,every node/.style={inner sep=0pt,font=\scriptsize}]
\path[use as bounding box] (0,0) rectangle (112,318);
\draw (17,8) -- (91,32);
\draw (2,30) .. controls (30,42) and (62,50) .. (90,55);
\draw (17,87) -- (89,20);
\draw (19,59) -- (92,115);
\draw (24,145) -- (73,80);
\draw (29,157) -- (75,83);
\draw (27,183) -- (89,220);
\draw (32,208) -- (98,262);
\draw (32,242) -- (79,185);
\draw (12,302) -- (79,233);
\draw (43,307) -- (98,241);
\node at (45,7) {$1$};
\node at (13,23) {$1$};
\node at (33,61) {$2$};
\node at (55,77) {$\unsure{2|1}$};
\node at (34,121) {$2$};
\node at (56,181) {$2$};
\node at (47,207) {$2$};
\node at (74,231) {$2$};
\node at (24,279) {$1$};
\node at (61,298) {$1$};
\fill (54,151) circle[radius=.45pt];
\fill (57,164) circle[radius=.45pt];
\fill (60,177) circle[radius=.45pt];
\end{tikzpicture}}} &
\parbox[c][5.45cm][c]{1.82cm}{\figure{%
\begin{tikzpicture}[x=.0158cm,y=-.0158cm,line width=.8pt,line cap=round,line join=round,every node/.style={inner sep=0pt,font=\scriptsize}]
\path[use as bounding box] (0,0) rectangle (112,318);
\draw (19,68) -- (75,13);
\draw (19,36) -- (78,95);
\draw (89,58) .. controls (43,124) and (32,192) .. (102,267);
\draw (18,116) -- (18,185);
\draw (12,158) -- (79,154);
\draw (32,298) -- (88,228);
\draw (23,248) -- (88,306);
\node at (49,16) {$1$};
\node at (57,67) {$2$};
\node at (29,126) {$1$};
\node at (73,146) {$2$};
\node at (68,197) {$3$};
\node at (49,253) {$2$};
\node at (81,293) {$1$};
\end{tikzpicture}}} &
\parbox[c][5.45cm][c]{1.82cm}{\figure{%
\begin{tikzpicture}[x=.0158cm,y=-.0158cm,line width=.8pt,line cap=round,line join=round,every node/.style={inner sep=0pt,font=\scriptsize}]
\path[use as bounding box] (0,0) rectangle (112,318);
\draw (37,38) -- (72,3);
\draw (38,17) -- (86,63);
\draw (30,99) -- (85,41);
\draw (21,74) .. controls (79,109) and (78,175) .. (14,207);
\draw (13,122) -- (101,120);
\draw (28,183) -- (89,250);
\draw (40,269) -- (86,208);
\draw (31,238) -- (92,289);
\node at (63,8) {$1$};
\node at (55,53) {$2$};
\node at (70,84) {$3$};
\node at (94,109) {$2$};
\node at (49,137) {$4$};
\node at (64,186) {$3$};
\node at (56,232) {$2$};
\node at (63,276) {$1$};
\end{tikzpicture}}} &
\parbox[c][5.45cm][c]{1.82cm}{\figure{%
\begin{tikzpicture}[x=.0158cm,y=-.0158cm,line width=.8pt,line cap=round,line join=round,every node/.style={inner sep=0pt,font=\scriptsize}]
\path[use as bounding box] (0,0) rectangle (112,318);
\draw (36,36) -- (67,0);
\draw (39,15) -- (79,57);
\draw (36,89) -- (77,41);
\draw (34,70) -- (83,119);
\draw (44,158) -- (88,100);
\draw (29,122) .. controls (44,147) and (81,159) .. (101,181) .. controls (110,191) and (111,207) .. (108,229);
\draw (30,192) .. controls (58,181) and (79,159) .. (95,133);
\draw (36,258) .. controls (64,235) and (88,231) .. (110,231);
\draw (45,242) .. controls (58,253) and (68,279) .. (64,296);
\node at (58,5) {$\unsure{1}$};
\node at (65,34) {$2$};
\node at (51,68) {$3$};
\node at (66,88) {$4$};
\node at (57,130) {$5$};
\node at (44,174) {$3$};
\node at (103,169) {$6$};
\node at (78,221) {$4$};
\node at (78,279) {$2$};
\end{tikzpicture}}} \\ \hline
\shortstack[l]{$\deleted{\unsure{n|m}}=\,$number of\\irred. components} &
1 & \nu & 1 & 2 & 3 & 5 & 5+\nu & 7 & 8 & 9 \\ \hline
\shortstack[l]{type of group $E(K)/E_0(K)$\\$\approx\widetilde X_0(k)/\widetilde E_0(k)$} &
(1)&({\unsure{\nu|v|2}})&(1)&(2)&(3)&(2\times2)&
\substack{(4),\,\nu\ \text{odd}\\(2\times2),\,\nu\ \text{even}}&
(3)&(2)&(1)\\ \hline
\widetilde E_0(k)\simeq E_0(K)/E_1(K) &
\widetilde E(k)&k^{*}&k^{+}&k^{+}&k^{+}&k^{+}&k^{+}&k^{+}&k^{+}&k^{+}\\ \hline
\multicolumn{11}{|@{}c@{}|}{%
\parbox[c]{250.3mm}{\setlength{\topsep}{0pt}\figure{%
\begin{tikzpicture}[x=.1726mm,y=.1726mm,line cap=round]
\path[use as bounding box] (0,0) rectangle (1450,78);
\foreach \divider in {278,395,512,629,746,863,980,1097,1214,1331}
  {\draw[line width=.4pt] (\divider,0) -- (\divider,78);}
\foreach \dx in {0,66,1330}{
  \begin{scope}[shift={(\dx,0)}]
    \draw[line width=1.5pt] (104,55) .. controls (68,55) and (61,43) .. (83,38)
      .. controls (114,31) and (98,20) .. (64,21);
  \end{scope}}
\node[anchor=west,inner sep=0pt] at (241,39)
  {\resizebox{184mm}{!}{BELOW THIS LINE THINGS ARE VALID ONLY FOR $\operatorname{CHAR}(k)\ne2,3$}};
\end{tikzpicture}}}}\\ \hline
v(\Delta_v)&0&\nu&2&3&4&6&6+\nu&8&9&10\\ \hline
\shortstack[l]{$v(\Delta_v)+1-n=f$\\$={}$exp. of $\mathfrak p$ in cond.}&
0&1&2&2&2&2&2&2&2&2\\ \hline
\text{behavior of }j&
v(j)\unsure{\geq|>}0&
\substack{\deleted{\unsure{v(j)|j}}\\v(j)=-\nu}&
\widetilde j=0&\widetilde j=1728&\widetilde j=0&
v(j)\geq0&v(j)=-\nu&\widetilde j=0&\widetilde j=1728&\widetilde j=0\\ \hline
\end{array}$}}

Here $\widetilde X_0$ denotes the non-singular part of the fiber. This is a (non-connected in general) algebraic group over $k$, whose connected component is $\widetilde E_0$. We have $E(K)/E_1(K)\approx\widetilde X_0(k)$, the isom induced by the reduction map, and assuming now $R$ complete. Note that if $p=\operatorname{char} k\ne2,3$, then we have minimality $\Longleftrightarrow$ \uline{either} $v(\Delta)<12$, \uline{or} $v(\Delta)+v(j)<12$. Also, if $p\ne2,3$, and $f=2$, then $E_0(K)$ is uniquely divisible by $2$ and $3$, while $E(K)/E_0(K)$ is killed by $12$, Hence $E(K)/E_0(K)$ is isomorphic to the group of points in $E(K)$ which are killed by $12$ in \unsure{this}%
\unsure{'|{}}
% archive sheet number: 9
\clearpage
\scanpage{11}\mspage{8.}
7. \uline{Algorithm} \uline{for} \uline{analysing} \uline{singular} \uline{fibers} (\uline{first} \uline{five} \uline{cases}). We assume now that our valuation ring $R$ is \uline{complete}, with \uline{perfect} residual field $k$. In connection with various conjectures, it is well to be able to compute effectively various invariants of an elliptic curve $E$ over $K$, to wit

the \uline{conductor} $f_E=\mathfrak p^{f}$, where $f=v(\Delta)+1-\unsure{n|m}$, $n$ being the number of components of $\widetilde X$ \uline{over} $\bar k$.

the group $E(K)/E_0(K)$\unsure{,|{}} whose order we denote by $c\deleted{\unsure{.|0}}$

the group $\widetilde E_0(k)$.

To compute these it is necessary to analyse the singular fiber $\widetilde X$ \unsure{à|a} la Neron, at least when $p=\operatorname{char} k\unsure{=|\ne}2\text{ or }3$. \reversemarginpar\margin{$\Vert$}[If $p=2$ or $3$, everything can be read off the table on page $6$, \deleted{\illegible{1}} if one notes the remarks at the bottom of page $6$; however the algorithm below applies in all cases.] When we refer to the ``\uline{type}'' we mean the type of the singular fiber $\widetilde X$ \uline{over} \uline{the} \uline{algebraic} \uline{closure} $\bar k$ \uline{of} $k$, which is designated by one of the $10$ Kodaira symbols.

To begin with, we simply assume an equation of the form (1.1) with coefficients $\boxed{a_i\in R}$; we do not assume it minimal. If it is not minimal, our algorithm will lead us to a minimalization of it. As we go along we make more and more assumptions. These are \uline{cumulative}, and are \fbox{boxed} for clarity. We include only brief remarks on proofs.

\[
\deleted{\unsure{0|1}}\;\text{\textcircled{1}}\quad
\mathfrak p\nmid\Delta\ \Longrightarrow\ \text{type }\mathrm I_0,\quad
 f=0,\quad c=1,\quad \widetilde E\ \text{\unsure{elliptic}}\,.
\]

\textcircled{2} Assume $\boxed{\mathfrak p\mid\Delta}$. then we can change coordinates so that $\boxed{\mathfrak p\mid a_3,a_4,\text{ and }a_6.}$ Do so. Then $\mathfrak p\nmid b_2\Longrightarrow\text{type }\mathrm I_{\unsure{\nu|v}}$, with $\nu=v(\Delta)$.

Let $k'$ be the splitting field over $k$ of the congruence $T^2+a_1T-a_2\equiv0$.
\[
\resizebox{\linewidth}{!}{$\displaystyle
\begin{aligned}
\text{\textcircled{\scriptsize 2a}}\;\;k'=k:\qquad&
 f=1,\quad c=v(\Delta),\quad \widetilde E_0(k)\approx k^{*}.\\[.6ex]
\text{\textcircled{\scriptsize 2b}}\;\;k'=k:\qquad&
 f=1,\quad c=\begin{cases}1,&\text{if }v(\Delta)\text{ odd}\\2,&\text{if }v(\Delta)\text{ even}\end{cases},\quad
 \widetilde E_0(k)=\parbox{.23\linewidth}{gp. of elts of $k'$ whose norm to $k$ is $1$.}
\end{aligned}$}
\]

\unsure{Proof}. This case \unsure{2)|\textcircled{2}} is the one in which $E$ can be described by $\theta$-functions, possibly twisted by an unramified extension. Everything clear from that point of view
\normalmarginpar
% archive sheet number: 10
\clearpage
\scanpage{12}\mspage{9.}
From now on $\widetilde E$ has a cusp and $\widetilde E_0(k)=k^+$

\textcircled{3} Assume $\boxed{\mathfrak p\mid b_2}$. Then $\mathfrak p^2\nmid a_6\Longrightarrow\text{type }\mathrm{II}$, and
\[
 f=v(\Delta),\qquad c=1,\qquad \widetilde E_0(k)=k^+.
\]
[Proof. Consider the $2$-dimensional local ring $A=R[x,y]_{\unsure{m|\mathfrak m}}$, where $m=(\mathfrak p,x,y)$. It is regular because $a_6\in(x,y)\Rightarrow m=(x,y)$. Hence Weierstrass model $=$ Neron model]

\textcircled{4} Assume $\boxed{\mathfrak p^2\mid a_6}$ (which implies $\mathfrak p^2\mid b_6$ and $b_8$). Then $\mathfrak p^3\nmid b_8\Rightarrow\text{Type }\mathrm{III}$, and
\[
 f=v(\Delta)-1,\qquad c=2,\qquad \widetilde E_0(k)=k^+
\]
\reversemarginpar\margin{\deleted{\unsure{\textcircled{4}|\textcircled{5}}}}[Proof. let $\mathfrak p=\pi R$, and let $a_i=\pi^m a_{i,m}$, $x=\pi^m x_m$, $y=\pi^m y_m$, etc. Our equation can be written
\[
 y_1^2+a_{1,0}x_1y_1\unsure{+|-}a_{3,1}y_1=\pi x_1^3+a_{2,0}x_1^2+a_{4,1}x_1+a_{6,2}
 \tag*{(7.1)}
\]
The singular point on the fiber (whose local ring was $A$) blows up into the conic
\[
 Y^2+a_{1,0}XY+a_{3,1}Y\equiv a_{2,0}X^2+a_{4,1}X+a_{6,2},
 \tag*{($*$)}
\]
whose descriminant is $b_8\pi^{-2}=b_{8,2}$.]

\textcircled{5} Assume $\boxed{\mathfrak p^3\mid b_8}$ (which implies $\mathfrak p^2\mid b_4$). Then $\mathfrak p^3\nmid b_6\Rightarrow\text{type }\mathrm{IV}$, and
\[
\resizebox{\linewidth}{!}{$\displaystyle\begin{aligned}
 f&=v(\Delta)-2,\quad
 c=\left\{
 \begin{aligned}
 &3\ \text{if}\quad T^2+a_{3,1}T-a_{6,2}\equiv0\quad\text{has roots in }k.\\
 &1\ \text{if}\qquad\qquad\quad \text{''}\qquad\qquad\quad\text{has roots }\notin k
 \end{aligned}\right.,\quad \widetilde E_0(k)=k^+
\end{aligned}$}
\]

[Proof: The conic ($*$) becomes $T^2+a_{3,1}T-a_{6,2}\unsure{=|\equiv}0$, where $T=Y-\alpha X$ is defined by $(Y-\alpha X)^2\equiv Y^2+a_1XY-a_2X^2$. The discriminant of $T^2+a_{3,1}T-a_{6,2}$ is $\pi^{-2} b_6=b_{6,2}$, so that if $\mathfrak p^3\nmid b_6$, then our conic degenerates into two \uline{distinct} lines. Dividing (7.1) by $x_1^3$, we get an equation of the form $F(U,V)=\pi$, where $F$ is a cubic with coeffs in $R$ which, mod $\pi$, factors into three distinct linear factors, $F\equiv L_1L_2L_3\pmod{\mathfrak p}$, such that the congruences $L_i\equiv0$ have a point in common. The local ring of that point is easily seen to be regular (the max ideal is generated by any two of the three factors $L_i$). Thus (7.1) gives a regular scheme over $R$ with fiber
\parbox[c]{1.48cm}{\setlength{\topsep}{0pt}\figure{\begin{tikzpicture}[x=1cm,y=1cm,line width=1.1pt,line cap=round]
\draw (-.49,-.07) -- (.49,.05);
\draw (-.34,.26) -- (.22,-.25);
\draw (-.36,-.30) -- (.21,.23);
\end{tikzpicture}}}
consisting of three lines $L_{\unsure{i|1}}=0$, $i\unsure{-|=}1,2,3$ meeting in a point. Concerning the value of $c$, we see that $c$ is equal to the number of these lines which are rational over $k$. One of them is, and the others are given by $T^2+a_{3,1}T-a_{6,2}\equiv0$.
\normalmarginpar
% archive sheet number: 11
\clearpage
\scanpage{13}\mspage{10}
8. \uline{Algorithm} \uline{continued}, \uline{the} \uline{last five} \uline{cases}. Assume now that $\boxed{\mathfrak p^3\mid b_6}$. Then we can change coördinates so that $\boxed{\mathfrak p\mid a_1\text{ and }a_2,\quad \mathfrak p^2\mid a_3\text{ and }a_4,\quad \text{and }\mathfrak p^3\mid a_6}$. This being done, consider the polynomial
\[
 P(T)=T^3+a_{2,1}T^2+a_{4,2}T+a_{6,3},
 \tag*{(8.1)}
\]
where the notation is as explained in the proof of step \textcircled{4}, page $9$. Our equation now becomes
\[
 \pi y_2^2+\pi a_{1,1}x_1y_2+\pi a_{3,2}y_2\unsure{-|=}P(x_1),
 \tag*{(8.2)}
\]
and our algorithm has three branches, according to the multiplicities of the roots of the congruence $P(T)\equiv0$:

\reversemarginpar
\margin{($1^{\mathrm{st}}$ branch)}\textcircled{6} If $P(T)\equiv0$ has three distinct roots, then we have Type $\mathrm I_0^{*}$, and
\[
 f=v(\Delta)-4,\quad\text{and}\quad c=1+(\text{number of roots of }P(T)\equiv0\text{ in }k)
\]

\margin{($2^{\mathrm d}$ branch)}\textcircled{7} If $P(T)\equiv0$ has one simple and one double root, then Type $\mathrm I_{\unsure{\nu|v}}^{*}$, $\nu>0$, and
\[
 f=v(\Delta)-4-\nu,\qquad c=2\text{ or }4,
\]
where $\nu$ and $c$ are obtained by the following procedure:

% margin note written sideways
\margin{\figure{\begin{tikzpicture}[x=1cm,y=1cm]
\node[rotate=90,inner sep=1pt] at (0,-3.1) {subprocedure branch \textcircled{7}};
\draw[line width=.8pt,line join=round] (1.25,0) -- (.22,0) -- (.22,-6.2);
\draw[line width=.8pt,line join=round] (.22,-5.94) -- (.30,-6.18) -- (.41,-5.96);
\end{tikzpicture}}}%
Translate $x$, so that $T\unsure{=|\equiv}0$ is a double root of $P(T)\equiv0$. Then $\mathfrak p^{\unsure{3|2}}\mid a_4$, $\mathfrak p^4\mid a_6$, $\mathfrak p^2\nmid a_2$ and (8.2) becomes
\[
 y_2^2+\pi a_{1,1}x_2y_2+a_{3,2}y_2
 =\pi^2x_2^3+\pi a_{2,1}x_2^2\unsure{+|-}\pi a_{4,\unsure{3|2}}x_2+a_{6,\unsure{4|1}}.
 \tag*{\fbox{7.1}}
\]

If $Y^2+a_{3,2}Y-a_{6,4}\equiv0$ has distinct roots, then
\[
 \nu=1,\quad\text{and}\quad c=\begin{cases}4&\text{if roots in }k.\\2&\text{otherwise }.\end{cases}
\]

If $Y^2+a_{3,2}Y-a_{6,\unsure{2|4}}\equiv0$ has a double root, we can translate $y$ so that the root is $Y\equiv0$. Then $\mathfrak p^3\mid a_3$, $\mathfrak p^5\mid a_6$, and our equation becomes
\[
 \pi y_3^2+\pi a_{1,1}x_2y_3+a_{3,3}y_3
 =\pi x_2^3+a_{2,1}{\unsure{x|Y|y}}_2^2+a_{4,3}x_2+a_{6,5}
 \tag*{\fbox{7.2}}
\]

If $a_{2,1}X^2+a_{4,3}X+a_{6,5}\equiv0$ has distinct roots, then
\[
 \nu=2,\ \text{and}\ c=\begin{cases}4&\text{if roots in }k\\2&\text{otherwise}.\end{cases}
\]
\normalmarginpar
\hfill(continued)
% archive sheet number: 12
\clearpage
\scanpage{14}\mspage{11.}
\reversemarginpar
% margin note written sideways
\margin{\figure{\begin{tikzpicture}[x=1cm,y=1cm]
\node[rotate=90,inner sep=1pt] at (0,-3.1) {subprocedure branch \textcircled{\unsure{7|1}}};
\draw[line width=.8pt,line join=round] (.10,-.22) -- (.17,0) -- (.51,-6.2) -- (1.75,-6.18);
\end{tikzpicture}}}%
If $a_{2,1}X^2+a_{4,3}X+a_{6,5}\equiv0$ has double root, then we can translate $x$ so that the root is $X\unsure{=|\equiv}0$, so that $\mathfrak p^4\mid a_4$ and $\mathfrak p^6\mid a_6$. Equation is now
\[
 y_3^2+\pi a_{1,1}x_3y_3+a_{3,3}y_3
 =\pi^3x_3^3\unsure{+|-}\pi^2a_{2,1}x_3^2+\pi^2a_{\unsure{1|4},4}x_3+a_{6,6}
 \tag*{\fbox{7.3}}
\]

If $Y^2+a_{3,3}\deleted{y}Y-a_{6,6}\equiv0$ has distinct roots, then
\[
 \unsure{\nu|v}=3,\quad\text{and}\quad c=\begin{cases}
 4&\text{if roots are in }k\\
 2&\text{if not}.
 \end{cases}
\]

If otherwise \unsure{....|...} etc. Keep going until the quadratic congruence which appears has distinct roots.... The process terminates, because the coefficients $a_3$, $a_4$, and $a_6$ are being made more and more highly divisible by $\mathfrak p$. Hence also $b_4$, $b_6$ and $b_8$, hence also $\Delta$. But $\Delta$ is fixed under all changes of coordinates (translations) involved. A crude estimate gives $\nu\leq\operatorname{ord}\Delta-3$ if I'm not mistaken.

\margin{(Branch 3 begins)}\textcircled{8} Suppose now $P(T)\equiv0$ has a triple root. We may assume the root is $T\equiv0$, so $\boxed{\mathfrak p^2\mid a_2,\ \mathfrak p^3\mid a_4,\ \text{and }\mathfrak p^4\mid a_6}$. The equation is
\[
 y_2^2+\pi a_{1,1}x_2y_2+a_{3,2}y_2
 =\pi^2x_2^3+\pi^2 a_{2,2}x_2^2+\pi a_{4,3}x_2+a_{6,4}
 \tag*{\fbox{8.1}}
\]
If $Y^2+a_{3,2}Y-a_{6,4}\equiv0$ has distinct roots, then Type $\mathrm{IV}^{*}$ and
\[
 f=v(\Delta)-6,\qquad c=\begin{cases}3&\text{if roots are in }k\\1&\text{if roots }\notin k.\end{cases}
\]

\margin{(Branch 3 continues)}\textcircled{9} Suppose $Y^2+a_{3,2}Y-a_{6,4}\equiv0$ has a double root. Then we can Assume it is $Y\unsure{=|\equiv}0$, so $\boxed{\mathfrak p^3\mid a_3,\ \mathfrak p^5\mid a_6}$. The equation is now
\[
 \pi y_3^2+\pi a_{1,1}x_2y_3+\pi a_{3,3}y_3
 =\pi x_2^3+\pi a_{2,2}x_2^{\unsure{2|7}}+a_{4,3}x_2+a_{6,5}
 \tag*{\fbox{9.1}}
\]
If $a_{4,3}\ne0$, i.e. if $\mathfrak p^4\nmid a_4$, then Type $\mathrm{III}^{*}$, and
\[
 f=v(\Delta)-7\qquad\text{and}\quad c=2
\]

\margin{(Branch 3 continues)}\textcircled{\scriptsize 10} Assume $\boxed{\mathfrak p^4\mid a_4}$. Then if $\mathfrak p^6\nmid a_6$ th\deleted{\unsure{a|e}}e type is $\mathrm{II}^{*}$ and
\[
 f=v(\Delta)-8\qquad\text{and}\quad c=1.
\]
\normalmarginpar

If $\mathfrak p^6\mid a_6$, original equation was not minimal. Start over with
\[
 y_3^2+a_{1,1}x_2y_3+a_{3,3}y_3=x_2^3+a_{2,2}x_2^2+a_{4,4}x_2+a_{6,6}\quad!
\]
% archive sheet number: 13
